\documentclass[a4paper,10pt,oneside,openany,final,article]{memoir}
\input{common}
\settocdepth{chapter}
\usepackage{minted}
\usepackage{fontspec}
% \setromanfont{Source Serif Pro}
% \setsansfont{Source Sans Pro}
% \setmonofont{Source Code Pro}

\begin{document}
\title{std::expected<T\&, E>}
\author{
  Steve Downey \small<\href{mailto:sdowney@gmail.com}{sdowney@gmail.com}> \\
}
\date{} %unused. Type date explicitly below.
\maketitle

\begin{flushright}
  \begin{tabular}{ll}
    Document \#: & D4280R0 \\
    Date: & \today \\
    Project: & Programming Language C++ \\
    Audience: & LEWG
  \end{tabular}
\end{flushright}

\begin{abstract}
  We propose to add a partial specialization of \tcode{std::expected} for
  lvalue reference types as the value type: \tcode{expected<T\&, E>}.
  The specialization holds a pointer to \tcode{T} when engaged and stores
  \tcode{E} by value when in error. Assignment has rebind semantics ---
  assigning to an engaged \tcode{expected<T\&, E>} always rebinds the
  reference, never assigns through to the referent. Construction from
  temporaries that would dangle is deleted. This mirrors the design of
  \tcode{optional<T\&>} accepted in \cite{P2988R12} and completes reference
  support across the standard library's principal sum types.
\end{abstract}

\tableofcontents*

\chapter*{Changes Since Last Version}

\begin{itemize}
\item \textbf{R0} Initial revision.
\end{itemize}

\chapter{Comparison table}

\section{Using a raw pointer for a fallible lookup}

The conventional idiom for a function that searches a container and may fail
is to return a raw pointer, or to take an out-parameter.  Neither approach
captures the error reason.

\begin{tabular}{ lr }
  \begin{minipage}[t]{0.45\columnwidth}
    \begin{minted}[fontsize=\small]{c++}
Cat* find_cat(std::string_view name);

Cat* cat = find_cat("Fido");
if (cat != nullptr) {
    return doit(*cat);
}
// error reason lost
    \end{minted}
  \end{minipage}
  &
    \begin{minipage}[t]{0.45\columnwidth}
      \begin{minted}[fontsize=\small]{c++}
auto find_cat(std::string_view name)
    -> std::expected<Cat&, Error>;

auto cat = find_cat("Fido");
return cat.and_then(doit);
      \end{minted}
    \end{minipage}
\end{tabular}

\section{Using \tcode{expected<Cat*, E>} as a substitute}

Returning \tcode{expected<Cat*, E>} is a common workaround.  It forces an
extra dereference through the pointer on the happy path, and the interface
still permits a null pointer to escape as a ``success''.

\begin{tabular}{ lr }
  \begin{minipage}[t]{0.45\columnwidth}
    \begin{minted}[fontsize=\small]{c++}
auto find_cat(std::string_view name)
    -> std::expected<Cat*, Error>;

auto cat = find_cat("Fido");
if (cat && *cat) {      // two checks
    return doit(**cat); // two derefs
}
    \end{minted}
  \end{minipage}
  &
    \begin{minipage}[t]{0.45\columnwidth}
      \begin{minted}[fontsize=\small]{c++}
auto find_cat(std::string_view name)
    -> std::expected<Cat&, Error>;

auto cat = find_cat("Fido");
return cat.and_then(doit);
      \end{minted}
    \end{minipage}
\end{tabular}

\section{Using \tcode{expected<reference\_wrapper<Cat>, E>}}

\tcode{std::reference_wrapper} avoids the null-pointer escape but is verbose
and requires \tcode{.get()} everywhere the reference is used.

\begin{tabular}{ lr }
  \begin{minipage}[t]{0.45\columnwidth}
    \begin{minted}[fontsize=\small]{c++}
auto find_cat(std::string_view name)
    -> std::expected<
           std::reference_wrapper<Cat>,
           Error>;

auto cat = find_cat("Fido");
if (cat) {
    return doit(cat->get());
}
    \end{minted}
  \end{minipage}
  &
    \begin{minipage}[t]{0.45\columnwidth}
      \begin{minted}[fontsize=\small]{c++}
auto find_cat(std::string_view name)
    -> std::expected<Cat&, Error>;

auto cat = find_cat("Fido");
return cat.and_then(doit);
      \end{minted}
    \end{minipage}
\end{tabular}

\section{Monadic composition}

The monadic API of \tcode{expected} works uniformly whether \tcode{T} is a
value type or a reference type, letting callers chain operations without
breaking out of the expected idiom.

\begin{tabular}{ lr }
  \begin{minipage}[t]{0.45\columnwidth}
    \begin{minted}[fontsize=\small]{c++}
Cat* cat = find_cat("Fido");
if (!cat) return make_error();
std::string* name = get_name(*cat);
if (!name) return make_error();
return process(*name);
    \end{minted}
  \end{minipage}
  &
    \begin{minipage}[t]{0.45\columnwidth}
      \begin{minted}[fontsize=\small]{c++}
return find_cat("Fido")
    .and_then(get_name)
    .and_then(process);
      \end{minted}
    \end{minipage}
\end{tabular}

\chapter{Motivation}

\tcode{std::expected<T, E>} \cite{P0323R12}, added in \CppXXIII{}, is a sum
type holding either a value of type \tcode{T} or an error of type \tcode{E}.
Like \tcode{std::optional} before it, the primary template explicitly
prohibits reference types for \tcode{T}; the current working draft contains
the static assertion:

\begin{minted}[fontsize=\small]{c++}
static_assert(!std::is_reference_v<T>,
    "T must not be a reference (use expected<T&,E> specialization)");
\end{minted}

The comment points toward the natural next step.

\tcode{std::optional<T\&>} has been accepted as \cite{P2988R12}.  Its
motivation section observes:

\begin{quote}
  Thus, we expect future papers to propose \tcode{std::expected<T\&,E>} and
  \tcode{std::variant} with the ability to hold references.
\end{quote}

This paper delivers the first of those two promises.

The design rationale for preferring rebind semantics over assign-through
semantics is well established by \cite{P2988R12} and the research surveyed
in \cite{P1683R0}.  All library implementations that have attempted
assign-through semantics for optional-like types over references have
eventually abandoned that approach due to the bug-prone behavior it
introduces.  The same applies to \tcode{expected<T\&, E>}.

Lookup functions that return references to elements of a container are
common in high-performance code precisely because they avoid copying
potentially large objects.  When such a lookup can fail, there is no
satisfying workaround available today:

\begin{itemize}
\item A raw pointer loses the error value and type safety.
\item \tcode{expected<T*, E>} requires two checks and two dereferences on
  the happy path, and does not prevent the caller from returning a null
  \tcode{T*} in the success channel.
\item \tcode{expected<std::reference_wrapper<T>, E>} works but is verbose;
  \tcode{.get()} is required everywhere the value is used, and
  \tcode{reference_wrapper} does not participate in the monadic API as
  naturally as a reference would.
\end{itemize}

\tcode{expected<T\&, E>} fills this gap cleanly.  It participates in the
monadic API of \tcode{expected} (\tcode{and_then}, \tcode{transform},
\tcode{or_else}, \tcode{transform_error}) without any awkward adapters.

There is a principled reason the standard chose to disallow reference types
for sum types: the semantics of assignment differ subtly between value types
and reference types.  For the primary template, \tcode{optional<T>} has pure
value semantics; assignment copies or moves the stored value.  For
\tcode{optional<T\&>}, assignment rebinds the pointer.  The observable
behavior is different.  However, this difference is not a defect --- it is
the same distinction that exists between \tcode{int} and \tcode{int*}, or
between \tcode{tuple<int>} and \tcode{tuple<int\&>}.  The standard already
accommodates reference types in \tcode{std::tuple}, \tcode{std::pair}, and
\tcode{std::reference_wrapper}.  Sum types should be no different.

The relationship between sum types and product types matters here.
\tcode{optional<T>} is equivalent to \tcode{variant<T, monostate>};
\tcode{expected<T, E>} is equivalent to \tcode{variant<T, E>}.  Neither
\tcode{variant} nor \tcode{tuple} can hold references today, but that is a
separate defect to address separately.  \tcode{expected<T\&, E>} can be
defined precisely and usefully now, following the exact same pattern as
\tcode{optional<T\&>}.

\chapter{Design}

The design of \tcode{expected<T\&, E>} follows \tcode{optional<T\&>}
\cite{P2988R12} closely.  Where the designs diverge it is because
\tcode{expected} carries an error value rather than simply being absent.

\section{Storage}

\tcode{expected<T\&, E>} stores a \tcode{T*} pointer when engaged and an
\tcode{E} value in-place when in the error state.  The default constructor
is deleted --- an engaged \tcode{expected<T\&, E>} must always hold a valid
reference, and there is no default reference.

\section{Rebind Semantics on Assignment}

Assignment always rebinds the pointer, never assigns through to the referent.
This is the only semantics that is not a source of latent bugs
\cite{P1683R0}.  Because an \tcode{expected<T\&, E>} in the value state
holds only a pointer, any assignment in the value-to-value case is equivalent
to pointer copy, and the full set of converting assignments can be
synthesized from copy-assignment via implicit construction of a temporary
\tcode{expected<T\&, E>} on the right-hand side.

\section{Dangling Prevention}

Constructors from types \tcode{U} that would bind \tcode{T\&} to a
temporary --- as defined by \tcode{reference_constructs_from_temporary} --- are
explicitly deleted.  This mirrors the approach taken for
\tcode{optional<T\&>} \cite{P2988R12}.

Deletion of these constructors can cause ambiguity in overload resolution
in the same way it does for \tcode{optional<T\&>}.  The design follows the
precedent set by \tcode{std::tuple} and \tcode{std::pair}: make dangling
calls ill-formed rather than silently eliding candidates.

\section{Value Access}

\tcode{operator*()} and \tcode{value()} return \tcode{T\&}.  Value category
is shallow: the reference stored in the \tcode{expected} is not propagated.
An \tcode{expected<T\&, E>\&\&} does not steal from the referent; it copies
the pointer.  This matches the behavior of \tcode{optional<T\&>}.

\section{Error Access}

\tcode{error()} returns a reference to the stored \tcode{E}, with the same
cv-ref qualifications as the primary template.

\section{value\_or}

\tcode{value_or} returns \tcode{T} by value, for the same reasons as in
\tcode{optional<T\&>::value_or} \cite{P2988R12}: returning a reference
would require the alternative to be of type \tcode{T\&} or compatible, which
is overly restrictive.  Free-function alternatives such as
\tcode{std::reference_or} \cite{P1255R12} and \tcode{std::reference_or}
\cite{D4270R0} provide reference-returning variants over all nullable types
for callers who need that behavior.

\section{Monadic Operations}

\tcode{and_then(f)}: calls \tcode{f} with \tcode{T\&}; returns the result of
\tcode{f}, which must be an \tcode{expected} specialization.

\tcode{transform(f)}: calls \tcode{f} with \tcode{T\&}; wraps the result in
an \tcode{expected}.

\tcode{or_else(f)}: calls \tcode{f} with the error \tcode{E}; returns the
result of \tcode{f}, which must be an \tcode{expected} specialization.

\tcode{transform_error(f)}: calls \tcode{f} with the error \tcode{E}; wraps
the new error in an \tcode{expected<T\&, G>}.

All four operations propagate the value or error, whichever is not being
processed, without modification.

\section{In-Place Construction}

\tcode{expected(std::in_place_t, Args...)} is deleted.  There is no place to
construct a \tcode{T} in an \tcode{expected<T\&, E>}.  In-place error
construction via \tcode{unexpect_t} works as in the primary template, because
the error type \tcode{E} is stored by value.

\section{Shallow vs Deep \tcode{const}}

\tcode{const expected<T\&, E>} permits modification of the referent through
\tcode{operator*()}.  Constness is shallow, consistent with the pointer model.
If deep constness is desired, use \tcode{expected<const T\&, E>}.

\section{Relationship to \tcode{expected<void, E>} and \tcode{expected<T, E>}}

The three specializations --- primary (\tcode{T} object type), void value
(\tcode{expected<void, E>}), and lvalue reference (\tcode{expected<T\&, E>})
--- share the same error-handling API.  Conversion from
\tcode{expected<U\&, G>} to \tcode{expected<T\&, E>} follows the same
constraints as conversion between \tcode{optional<U\&>} and
\tcode{optional<T\&>}.

\section{Feature Test Macro}

A feature test macro \tcode{__cpp_lib_expected_ref} is proposed, allowing
code to detect the availability of the reference specializations.

\chapter{Proposal}

Add a partial specialization \tcode{expected<T\&, E>} to the \tcode{<expected>}
header, with the design described above.

The reference implementation is available as part of the Beman Project's
\tcode{beman.expected26} library \cite{Downey_beman_expected}.

\chapter{Impact on the standard}

A pure library extension.  No changes to the language are required.  No
existing well-formed code is made ill-formed.

The change adds a new partial specialization of \tcode{expected}; code that
currently relies on the absence of such a specialization (e.g., by static
assertion) will need to be updated, but such code would necessarily have been
deliberately excluding a feature rather than relying on normal program
behavior.

\chapter{Wording}

Wording to follow in a subsequent revision.

\chapter{Acknowledgments}

The design of \tcode{expected<T\&, E>} is a direct extension of the work on
\tcode{optional<T\&>} \cite{P2988R12} by Steve Downey and Peter Sommerlad.
Many thanks to JeanHeyd Meneide for the foundational research in
\cite{P1683R0} and \cite{REFBIND}.

Thanks to the contributors to the Beman Project's \tcode{beman.expected26}
reference implementation \cite{Downey_beman_expected}.

\chapter*{Document history}

\begin{itemize}
\item \textbf{R0} Initial revision.
\end{itemize}

\renewcommand{\bibname}{References}
\bibliographystyle{abstract}
\bibliography{wg21,mybiblio}

\backmatter
\chapter*{Implementation}

\inputminted{c++}{../include/beman/expected/expected.hpp}

\end{document}
